\section{Related Works}

The study of using optimal transport \cite{Villani2008OptimalTO} to improve flow matching was proposed firstly in \cite{tong2024improvinggeneralizingflowbasedgenerative, pooladian2023multisampleflowmatchingstraightening}, which adds optimal transport induced coupling for the minibatch before the interpolation and model training step. Several follow-up works \cite{kornilov2024optimalflowmatchinglearning, yue2025oatfmoptimalaccelerationtransport} have aimed to further expand the framework. All have assumed that the sample space is $\mathbb{R}^d$. In order to better capture the resolution-invariance of grid data and continuity of path data, Functional Flow Matching \cite{kerrigan2023functionalflowmatching} first proposed a framework to do flow matching in Hilbert space, relying on Gaussian measure theory. However, the same optimal transport acceleration has not been made in this space.

The line of research that embeds probability measures in RKHS has been proposed by the pioneering work of \cite{Sriperumbudur2010Hilbert, gretton2008kernelmethodtwosampleproblem}, where the kernel-induced metrics, called Intergral Probability Metrics (IPM), has been studied extensively \cite{Dudley_2002, Sriperumbudur2010Hilbert} and used for two-sample  hypothesis testing \cite{gretton2008kernelmethodtwosampleproblem}, most recently in path space \cite{chevyrev2022signaturemoments}. IPW in the form of Maximum Mean Discrepancy (MMD) has been linked with both probability divergence and Wasserstein distance \cite{feydy19a-pmlr} and has seen clear entropic optimal transport formulation in \cite{Li_2021_CVPR} in the form of Hilbert Sinkhorn Divergence and a gradient flow formulation in \cite{galashov2024deepmmdgradientflow, arbel2019maximummeandiscrepancygradient}. Nevertheless, the extension of these frameworks to function space is an active area of research in mathematics, and our work presents the first to explore it in the context of flow matching. 